% Standalone problem document, generated from the adjacent README.md.
% Compile with XeLaTeX. Mathematical content is maintained in Markdown.
\documentclass[11pt,a4paper]{article}
\usepackage[fontsize=10.5pt]{scrextend}
\usepackage[margin=22mm,headheight=15pt,headsep=9mm,footskip=11mm]{geometry}
\usepackage{fontspec}
\setmainfont{texgyrepagella}[Extension=.otf,UprightFont=*-regular,BoldFont=*-bold,ItalicFont=*-italic,BoldItalicFont=*-bolditalic]
\setsansfont{texgyreheros}[Extension=.otf,UprightFont=*-regular,BoldFont=*-bold,ItalicFont=*-italic,BoldItalicFont=*-bolditalic]
\setmonofont{texgyrecursor}[Extension=.otf,UprightFont=*-regular,BoldFont=*-bold,ItalicFont=*-italic,BoldItalicFont=*-bolditalic,Scale=MatchLowercase]
\usepackage{amsmath,amssymb}
\usepackage{unicode-math}
\setmathfont{texgyrepagella-math.otf}
\usepackage{xcolor,microtype,enumitem,needspace,fancyhdr}
\usepackage{bookmark}
\definecolor{ink}{HTML}{17344B}
\definecolor{accent}{HTML}{197D80}
\definecolor{muted}{HTML}{596773}
\hypersetup{colorlinks=true,linkcolor=accent,urlcolor=accent,pdftitle={SP-10 - The
graph complement conjecture for minimum symmetric
rank},pdfauthor={Open Problems in Numerical Linear Algebra}}
\urlstyle{same}
\setlength{\parindent}{0pt}
\setlength{\parskip}{5pt plus 1pt minus 1pt}
\linespread{1.04}
\setlength{\emergencystretch}{3em}
\widowpenalty=10000
\clubpenalty=10000
\displaywidowpenalty=10000
\allowdisplaybreaks[1]
\setlist{nosep,leftmargin=1.5em}
\providecommand{\tightlist}{\setlength{\itemsep}{0pt}\setlength{\parskip}{0pt}}
\providecommand{\pandocbounded}[1]{#1}
\pagestyle{fancy}
\fancyhf{}
\fancyhead[L]{\sffamily\footnotesize\color{muted}OPEN PROBLEMS IN NUMERICAL LINEAR ALGEBRA}
\fancyhead[R]{\sffamily\bfseries\footnotesize\color{accent}SP-10}
\fancyfoot[L]{\parbox[b]{0.9\textwidth}{\sffamily\fontsize{8}{10}\selectfont\color{muted}Editorial ratings; a bounded search cannot certify that a problem remains unsolved.\\Literature check: 2026-09-10}}
\fancyfoot[R]{\sffamily\footnotesize\color{muted}\thepage}
\renewcommand{\headrulewidth}{0.3pt}
\renewcommand{\headrule}{\hbox to\headwidth{\color{accent}\leaders\hrule height \headrulewidth\hfill}}
\makeatletter
\renewcommand{\section}{\Needspace{5\baselineskip}\@startsection{section}{1}{0pt}{12pt plus 2pt minus 2pt}{4pt}{\sffamily\bfseries\large\color{ink}}}
\renewcommand{\subsection}{\Needspace{4\baselineskip}\@startsection{subsection}{2}{0pt}{9pt plus 1pt minus 1pt}{3pt}{\sffamily\bfseries\normalsize\color{ink}}}
\makeatother
\setcounter{secnumdepth}{0}
\begin{document}
{\sffamily\small\color{muted}Eigenvalues and inverse problems\par}
\vspace{3pt}
{\raggedright\hyphenpenalty=10000\exhyphenpenalty=10000\sffamily\bfseries\fontsize{21}{24}\selectfont\color{ink}The
graph complement conjecture for minimum symmetric rank\par}
\vspace{7pt}
{\sffamily\small\textbf{Difficulty:} extreme\quad\textbf{Importance:} interesting
to the community\par}
{\sffamily\small\bfseries\color{accent}Status: Partially resolved\par}
\vspace{4pt}
{\color{accent}\hrule height 0.8pt}
\vspace{6pt}
\textbf{Rating rationale:} Extreme reflects a longstanding bound over
all symmetric sparsity patterns, still beyond the available general rank
estimates. Community impact comes from its central role in graph inverse
eigenvalue problems and complementary low-rank realizations.

\subsection{Problem statement}\label{problem-statement}

For a finite simple undirected graph \(G\) with vertex set
\(\{1,\ldots,n\}\), \(n\ge1\), define \[
\mathcal S(G)=\{A\in\mathbb R^{n\times n}:A=A^T,\ \mathcal G(A)=G\},
\qquad
\operatorname{mr}(G)=\min_{A\in\mathcal S(G)}\operatorname{rank}A.
\] Here \(\mathcal G(A)\) has an edge \(\{i,j\}\) exactly when
\(i\ne j\) and \(a_{ij}\ne0\). Diagonal entries are unrestricted. Let
\(G^c\) have exactly the complementary edges between distinct vertices.
Does every such graph satisfy \[
\operatorname{mr}(G)+\operatorname{mr}(G^c)\le n+2?
\]

\subsection{Relevance and ratings}\label{relevance-and-ratings}

This is a structured low-rank feasibility problem for complementary
symmetric sparsity patterns. Because diagonal shifts preserve those
patterns, minimum rank also determines the largest eigenvalue
multiplicity realizable with a prescribed off-diagonal pattern. It is an
inverse eigenvalue problem, not a claim about the rank of the unweighted
adjacency matrix.

\subsection{References}\label{references}

\begin{itemize}
\item
  F. Barioli, S. M. Fallat, H. Gupta, and Z. Li, \emph{The weak version
  of the graph complement conjecture and partial results for the delta
  conjecture}, Discrete Mathematics 349 (2026), 114861, Conjecture 1.1
  (\href{https://arxiv.org/html/2505.24577v1}{primary manuscript};
  \href{https://doi.org/10.1016/j.disc.2025.114861}{journal}).
\item
  S. M. Fallat and L. Hogben, \emph{The minimum rank of symmetric
  matrices described by a graph: A survey}, Linear Algebra and its
  Applications 426 (2007), 558--582, introduction and discussion of
  minimum rank
  (\href{https://aimath.org/WWN/matrixspectrum/FallatHogbenMinRank07.pdf}{primary
  manuscript};
  \href{https://doi.org/10.1016/j.laa.2007.05.036}{journal}).
\item
  M. Jansrang and S. K. Narayan, \emph{Graph complement conjecture for
  classes of shadow graphs}, Operators and Matrices 15 (2021), 589--614,
  §1 and Theorem 1
  (\href{https://files.ele-math.com/articles/oam-15-40.pdf}{primary
  paper}; \href{https://doi.org/10.7153/oam-2021-15-40}{journal}).
\end{itemize}

\subsection{Status check}\label{status-check}

The 2026 paper resolves a weak version with a coefficient strictly below
\(2\) in front of \(n\), not the displayed coefficient \(1\). Searches
for ``graph complement conjecture'', ``minimum rank'', ``solved'',
``proof'', and 2025--2026 found no general resolution. The source's May
2025 arXiv version and the 2026 journal record were checked. No
positive-semidefinite or strong-Arnold variant is separately counted
here.

\subsection{Audit update --- 2026-09-10}\label{audit-update-2026-09-10}

\textbf{Partially resolved:} trees, unicyclic graphs, chordal graphs and
partial 3-trees satisfy the stronger real positive-semidefinite version;
Jansrang--Narayan §1 records these results and Theorem 1 proves
additional shadow-graph families. Since unrestricted minimum rank is no
larger than PSD minimum rank, these are cases of the displayed target.
The general graph case remains the unresolved question. The full primary
manuscript and 2026 publication record above, the 2021 paper, and
targeted later-resolution searches were checked; no full resolution was
located.
\end{document}
